\section{Zadanie 7}
\autor{Mateusz Wojaczek}
Dana jest moneta asymetryczna z nieznanym Ci \( p \in (0, 1) \).
Jak wiele rzutów musisz wykonać, aby z prawdopodobieństwem co najmniej \( 0.9 \) poznać wartość \( p \) z dokładnością do \( 0.1 \)?

\subsection*{Rozwiązanie}
\[
\text{Niech:}
\]
\[
\text{orzeł } \rightarrow p, \quad \text{reszka } \rightarrow 1 - p
\]

\[
X - \text{ liczba wyrzuconych orłów w } n \text{ rzutów}
\]

\[
E[X] = np
\]

\[
P((p+0{,}1)n \geq X \geq (p-0{,}1)n) \geq 0{,}9 
\]
\[
\Leftrightarrow P(|X - pn| \geq 0{,}1n) \leq 0{,}1
\]

\[
\mathrm{Var}(X) = \sum_{i=1}^{n} \mathrm{Var}(X_i) = n \big(E[X_i^2] - E[X_i]^2\big) = n(p - p^2) = np(1 - p)
\]

\[
P(|X - pn| \geq 0{,}1n) = P(|X - E[X]| \geq 0{,}1n) \leq \frac{\mathrm{Var}(X)}{(0{,}1n)^2}
\]
\[
= \frac{np(1 - p)}{n^2} \cdot 100 = \frac{p(1 - p)}{n} \cdot 100
\]

\[
\frac{p(1 - p)}{n} \cdot 100 \leq 0{,}1
\]
\[
p(1 - p) \cdot 1000 \leq n
\]

Dla \( p = \frac{1}{2} \) lewa strona osiąga wartość największą, zatem:
\[
\frac{1}{4} \cdot 1000 \leq n
\]


\subsection*{Odpowiedź}
\[
\boxed{n \geq 250}
\]
